\documentclass[11pt]{article}
\usepackage[margin=1in]{geometry}
\usepackage{amsmath}
\usepackage{amssymb}
\usepackage{booktabs}
\usepackage{enumitem}
\usepackage{hyperref}

\title{Scalability and Code Review Notes for the Current GQE Methodology}
\author{}
\date{}

\begin{document}
\maketitle

\section{Scope}

This note reviews the current repository without changing implementation code.  The
review focuses on the methodology currently implemented in
\texttt{main.py}, \texttt{trainer.py}, \texttt{policy.py},
\texttt{circuit.py}, \texttt{refine.py}, \texttt{pareto.py}, and
\texttt{simplify.py}; it also checks benchmark/reporting paths that affect how
results are interpreted.

The existing test suite was run with \texttt{pytest -q}.  It passed:
\[
  30 \text{ passed.}
\]
The test run emitted JAX persistent-cache warnings because the configured cache
directory under \texttt{\$HOME/.cache/jax-compilation-cache} was read-only in the
current sandbox.  The warnings do not invalidate the tests, but they matter for
long repeated runs.

\section{Current Methodology}

The current implementation is a target-conditioned, hybrid-action GQE search:

\begin{enumerate}[leftmargin=*]
  \item \textbf{Gate vocabulary.}  \texttt{operator\_pool.py} builds one token
  for each configured one-qubit rotation on each qubit, one \texttt{SX} token per
  qubit, and every directed \texttt{CNOT} pair.  With \(n\) qubits and \(r\)
  configured rotation axes, the gate pool has
  \[
    n(r+1) + n(n-1)
  \]
  gate tokens, plus \texttt{BOS}, \texttt{STOP}, and \texttt{PAD}.

  \item \textbf{Target conditioning.}  The target unitary is flattened into
  real and imaginary parts and projected into the transformer embedding space.
  This vector has length \(2 \cdot 4^n\).  For 4 qubits this is still small
  (\(512\) floats), but it grows exponentially.

  \item \textbf{Hybrid policy.}  \texttt{HybridPolicy} is a causal transformer.
  At every sequence position it emits discrete token logits and a Gaussian
  angle head \((\mu,\log\sigma)\).  Rollout sampling uses a cached autoregressive
  decoder in \texttt{build\_rollout\_fn}.  Discrete sampling is from
  \(\mathrm{softmax}(-\beta \,\mathrm{logits})\), so the logits are being treated
  like energies rather than conventional larger-is-better logits.

  \item \textbf{Circuit evaluation.}  \texttt{CircuitEvaluator} composes the
  full \(2^n \times 2^n\) unitary by applying each gate to the rows of the
  running matrix.  Fidelity is process fidelity,
  \[
    F(U,V) = \frac{|\mathrm{Tr}(U^\dagger V)|^2}{d^2},
    \qquad d = 2^n,
  \]
  which is invariant to global phase.

  \item \textbf{Angle refinement.}  Rollouts may be scored after an inner Adam
  refinement of all sampled angles.  A final archive refinement also runs after
  training, optionally followed by host-side sweep polishing.

  \item \textbf{Reward.}  The reward is lexicographic in spirit.  Fidelity is
  optimized first using a log-infidelity reward.  CNOT and depth penalties only
  activate when fidelity exceeds \texttt{lex\_structure\_fidelity\_threshold}.
  A sequence discount and no-\texttt{STOP} penalty encourage shorter, terminated
  rows.

  \item \textbf{Policy update.}  The discrete token policy is updated with a
  PPO/GRPO-style sequence-level objective.  The stored old log probabilities are
  only the discrete-token log probabilities.  The continuous angle head is
  trained by supervised regression toward the stored/refined angles, not by a
  reward-weighted continuous-action policy-gradient term.

  \item \textbf{Pareto archive.}  Candidate circuits above a fidelity floor are
  archived under four objectives: maximize fidelity, minimize depth, minimize
  total gates, and minimize CNOT count.  After final refinement, the archive is
  simplified algebraically and rebuilt.

  \item \textbf{Simplification.}  \texttt{simplify.py} applies local rewrite
  rules: same-axis rotation merging, identical CNOT cancellation, zero-angle
  removal, angle reduction, partial \texttt{SX} periodicity, and two CNOT
  commutation look-through rules (\(R_Z\) on the CNOT control and \(R_X\) on the
  CNOT target).
\end{enumerate}

\section{Interpretation of the 2-Qubit Improvement}

The strong 2-qubit results after simplification are plausible because the policy
often samples redundant circuits.  Post-hoc algebraic simplification can remove
many redundant rotations and CNOT pairs without changing process fidelity.  The
current benchmark logs show 2-qubit brickwork runs reaching \(F=1.0\) with
competitive or better CNOT counts after simplification.

The important caveat is that simplification is currently mostly a reporting and
archive post-processing step.  The policy is trained and rewarded using the
pre-simplification structure during rollouts.  Therefore, the improvement shows
that the sampled circuits contain useful solutions hidden behind redundant
syntax, but it does not yet mean the policy has learned to sample compact
canonical circuits directly.  For 4+ qubits, this distinction becomes much more
important because redundant syntax grows rapidly with sequence length.

\section{Main Scalability Risks}

\begin{table}[h]
\centering
\begin{tabular}{@{}lll@{}}
\toprule
Component & Current scaling & Practical issue for 4+ qubits \\
\midrule
Dense unitary evaluation & \(O(L4^n)\) per circuit & Expensive with long horizons and refinement \\
Target feature vector & \(O(4^n)\) & Fine at 4q, increasingly large after that \\
Transformer PPO forward & \(O(BT^2)\) & Very costly when \(T\) reaches hundreds \\
Rollout decoder & roughly \(O(BT^2)\) & KV cache still attends over full length each step \\
Inner refinement & \(O(BS_{\rm ref}L4^n)\) & Dominant cost if applied to every sample \\
Pareto update batch & \(O((A+K)^2)\) & Manageable at current sizes, risky if \(K\) grows \\
Qiskit dense reference & exponential & Can dominate or fail for 4+ dense targets \\
\bottomrule
\end{tabular}
\end{table}

The current default action horizon is especially important.  The function
\texttt{rollout\_context\_tokens} estimates an action budget from a generic CNOT
lower-bound-like expression:
\[
  \left\lceil \frac{4^n - 3n - 1}{4} \right\rceil.
\]
It then uses roughly \(7.5\) times that value, with a minimum of \(32\) and a
maximum of \(1023\) actions.  This gives:
\[
  T_{\rm actions}(2q)=32,\quad
  T_{\rm actions}(3q)\approx105,\quad
  T_{\rm actions}(4q)\approx458,\quad
  T_{\rm actions}(5q)=1023 \text{ (cap)}.
\]
For structured brickwork targets, these horizons can be far larger than needed.
For dense Haar targets, they may still be insufficient or inefficient because the
search space is enormous.

\section{Recommendations for Scaling to 4+ Qubits}

\subsection{Make Simplification Part of the Search Loop}

The biggest methodological change should be to score and archive simplified
circuits during rollout, not only after training.  A recommended pipeline is:

\begin{enumerate}[leftmargin=*]
  \item Decode sampled tokens and angles.
  \item Simplify/canonicalize the circuit.
  \item Recompute simplified depth, total gates, and CNOT count.
  \item Use the simplified structural metrics in the reward.
  \item Store the simplified sequence in the archive using a canonical hash.
\end{enumerate}

This gives the policy credit for circuits that simplify well and prevents the
archive from filling with syntactically different versions of the same circuit.
For 4+ qubits, this is likely more valuable than simply increasing model size.

\subsection{Use Target-Aware Action Budgets}

The action budget should be configurable and target-aware.  For 4-qubit
brickwork targets, a horizon of \(458\) is probably too expensive and too
permissive.  Start with shorter horizons such as \(64\), \(96\), or \(128\),
then increase only if the Pareto front saturates below the desired fidelity.

For dense 4-qubit Haar targets, use a separate configuration because the required
gate count is much larger.  Do not use the same horizon and sampling budget for
brickwork, random reachable, and full Haar targets.

\subsection{Refine Only Promising Skeletons}

The inner refiner is very expensive because it optimizes every sampled circuit.
For 4+ qubits, use a two-stage scoring scheme:

\begin{enumerate}[leftmargin=*]
  \item Evaluate raw fidelities for all samples.
  \item Select the top \(K\) candidates by raw fidelity, archive novelty, or
  simplified compactness.
  \item Run Adam refinement only on those \(K\) candidates.
  \item Use exact refined scores for archive updates.
\end{enumerate}

The angle head can still learn from the refined candidates.  This keeps the
benefit of refinement without multiplying the full rollout cost by the number of
Adam steps for every sample.

\subsection{Constrain the CNOT Topology When Appropriate}

The current pool includes all directed CNOTs.  For \(n=4\), this is \(12\) CNOT
tokens; for larger \(n\), it grows as \(n(n-1)\).  If the target is brickwork or
hardware-local, restrict CNOTs to the relevant graph.  This reduces the token
space, removes many bad skeletons, and makes Qiskit comparisons more meaningful.

A good next experiment is a 4-qubit nearest-neighbor brickwork setup with only
adjacent CNOT directions, then compare against an all-to-all run.

\subsection{Use Hierarchical or Blockwise Synthesis}

For 4+ qubits, especially dense targets, direct monolithic sequence search is
not the most scalable route.  Better options are:

\begin{itemize}[leftmargin=*]
  \item decompose the target into two-qubit or nearest-neighbor blocks, then use
  GQE to optimize each block or boundary;
  \item use known decompositions such as KAK for two-qubit blocks and CSD/QSD for
  larger unitaries, then let GQE reduce or rewrite the resulting circuit;
  \item train the policy to synthesize and simplify local windows rather than
  entire dense unitaries from scratch.
\end{itemize}

This also matches the observation that simplification gives large gains: a local
rewrite/synthesis model can specialize in removing redundancies introduced by a
baseline decomposition.

\subsection{Use Approximate Screening Before Exact Unitary Fidelity}

Dense exact fidelity is still feasible at 4 qubits, but the combination of long
horizons, many samples, and refinement is expensive.  For larger systems, screen
candidates with cheaper metrics first:

\begin{itemize}[leftmargin=*]
  \item random input-state probes;
  \item selected Pauli-transfer or Hilbert-Schmidt sketches;
  \item target-specific local observables for structured circuits.
\end{itemize}

Only top candidates should receive full \(2^n \times 2^n\) unitary evaluation and
angle refinement.

\subsection{Use Curriculum and Warm Starts}

The current runs initialize a fresh policy for each target.  For 4+ qubits, use
curriculum:

\begin{itemize}[leftmargin=*]
  \item train on reachable/random shallow targets before harder brickwork or
  Haar-like targets;
  \item increase qubit count, target depth, and horizon gradually;
  \item pretrain on many target circuits so the target-conditioned policy learns
  reusable structure;
  \item warm-start 4-qubit runs from 3-qubit or smaller-depth checkpoints.
\end{itemize}

\section{Confirmed Bugs and Implementation Issues}

\subsection{High Priority: Simplified Angle Alignment Breaks Reported Qiskit Circuits}

Files involved: \texttt{simplify.py} and \texttt{reporting.py}.

\texttt{simplify.py:\_encode} returns an angle vector with a leading BOS-aligned
zero:
\[
  [0,\theta_1,\theta_2,\ldots].
\]
However \texttt{reporting.py:build\_qiskit\_circuit\_from\_actions} removes BOS
from the token sequence and then reads \texttt{angles[pos]}.  That function
expects action-aligned angles:
\[
  [\theta_1,\theta_2,\ldots].
\]

I confirmed this with a one-gate 2-qubit probe.  A simplified circuit containing
\(R_Z(0.7)\) stored angles as
\[
  [0, 0.699999988, 0, 0].
\]
When passed to reporting, the Qiskit circuit used angle \(0\) for the first gate.
The archive-reported fidelity was \(1.0\), while the Qiskit-recomputed fidelity
was \(0.8824210975\).

Impact: benchmark Pareto metrics can still be mathematically meaningful if the
simplification itself is correct, but exported/reported Qiskit circuits after
simplification can be wrong.  This can make \texttt{verified\_fidelity} and
\texttt{qc\_fidelity} disagree.  Before using simplified circuits as final
artifacts, the project needs one consistent angle convention.

\subsection{High Priority: Documented SX Periodicity Is Not Fully Implemented}

Files involved: \texttt{simplify.py} and \texttt{docs/simplification\_rules.tex}.

The documentation says:
\[
  \mathrm{SX}^2 = R_X(\pi), \qquad \mathrm{SX}^4 = I.
\]
The implementation only rewrites \(\mathrm{SX}^2\) when an \texttt{RX} token
exists in the configured pool.  With the current default-style pool
\texttt{["rz", "ry"]}, four consecutive \texttt{SX} gates remain unchanged.

With an \texttt{RX}-enabled pool, four \texttt{SX} gates simplify only to
\(R_X(2\pi)\), not to identity.  This is because the code reduces modulo
\(4\pi\) and only drops angles close to \(0\), while the documentation states
that \(2\pi\) rotations may be dropped as global phase under the process-fidelity
objective.

Impact: the simplifier leaves avoidable gates in exactly the kind of redundant
circuits that become common in long 4+ qubit rollouts.

\subsection{High Priority: Simplification Is Not Used in the Rollout Reward}

Files involved: \texttt{trainer.py} and \texttt{simplify.py}.

During rollout, CNOT count, depth, total gates, reward, and archive admission are
computed from the raw sampled sequence.  Simplification is applied only after
training to the final archive.  This means compact-after-rewrite circuits do not
receive structural reward credit during learning.

Impact: for 4+ qubits, the policy may learn to rely on post-processing instead
of learning compact canonical skeletons.  The larger the horizon, the worse this
credit-assignment mismatch becomes.

\subsection{High Priority: One-Qubit Evaluation Is Broken Although Config Allows It}

Files involved: \texttt{config.py} and \texttt{circuit.py}.

\texttt{validate\_config} allows \texttt{num\_qubits=1}.  But
\texttt{CircuitEvaluator.\_build\_unitary} always computes both the one-qubit and
two-qubit branches before selecting with \texttt{jnp.where}.  For \(n=1\),
\(\texttt{d // 4}=0\), so the two-qubit branch has incompatible row-block shapes.
A direct 1-qubit fidelity call fails with:
\[
  \texttt{TypeError: add got incompatible shapes for broadcasting: (0, 2), (2, 2).}
\]

Impact: not directly a 4-qubit blocker, but it is a correctness gap and a sign
that branch evaluation in the JAX circuit evaluator should be tested more
thoroughly.

\subsection{Medium Priority: Continuous Angle Log Probabilities Are Not Used by PPO}

Files involved: \texttt{trainer.py}, \texttt{policy.py}, and \texttt{loss.py}.

The rollout path samples continuous angles and there is code for wrapped-normal
log probabilities, but the PPO update stores and uses only discrete-token
\texttt{log\_p\_disc}.  \texttt{log\_p\_cont} is computed in one helper but is not
part of the sequence PPO ratio.  The angle head is trained with a supervised
\(1-\cos(\Delta\theta)\) loss toward stored refined angles.

This may be intentional, and it can work when angle refinement supplies good
targets.  The limitation is that reward does not directly reinforce angle
sampling.  As qubit count and angle dimension grow, this may make the policy
depend more heavily on expensive classical refinement.

\subsection{Medium Priority: No Fidelity Reverification After Simplification}

Files involved: \texttt{simplify.py}.

\texttt{simplify\_point} copies the old fidelity into the new \texttt{ParetoPoint}
without reevaluating the simplified sequence.  The rewrite rules are intended to
be exact, but the confirmed angle-alignment issue shows why a verification pass
is important.

Recommendation: after simplification, recompute fidelity for at least debug mode
or archive-finalization mode and assert that the process fidelity did not drop
beyond a small tolerance.

\subsection{Medium Priority: Archive Can Store Duplicate Objective Points}

Files involved: \texttt{pareto.py}.

Dominance requires strict improvement in at least one objective.  Two different
circuits with identical \((F,\mathrm{depth},\mathrm{gates},\mathrm{CNOTs})\) do
not dominate each other, so duplicates can remain until crowding pruning removes
something.  For 4+ qubits, this wastes archive capacity.

Recommendation: store a canonical simplified circuit hash and reject exact
duplicates, or keep only the shorter/canonical representative for each objective
tuple.

\subsection{Medium Priority: Action Horizon Is Hard-Coded Rather Than Configurable}

Files involved: \texttt{trainer.py}.

\texttt{rollout\_context\_tokens} is not exposed in \texttt{config.yml}.  Because
it grows quickly with \(n\), 4-qubit runs can silently move to hundreds of
actions.  This affects transformer memory, rollout time, refinement time, and
archive complexity.

Recommendation: make the action horizon a config field with sensible
target-type defaults.

\subsection{Medium Priority: Benchmark Outputs Are Not Crash-Resilient}

Files involved: \texttt{benchmark.py}.

\texttt{benchmark.py} contains helper functions for JSONL loading, completed
seed detection, summary writing, and metadata writing, but the main function does
not currently write per-seed JSONL rows or the summary/metadata files.  It keeps
rows in memory and writes final CSV/plot/log outputs at the end.

Impact: long 4-qubit or 5-qubit benchmarks can lose most results if interrupted.
This is a scalability and reproducibility problem.

\subsection{Medium Priority: Run Artifacts Do Not Store Reconstructable Circuits}

Files involved: \texttt{reporting.py}.

\texttt{save\_run\_artifact} stores target matrices, config, logs, and Pareto
metrics, but not the token sequences or optimized angle arrays for Pareto
entries.  This makes it hard to verify, rerender, or post-process the exact
circuits later.

Recommendation: store token sequences, angle arrays, and the angle convention
explicitly.

\subsection{Lower Priority: Hypervolume Is Not Comparable Across Runs}

Files involved: \texttt{pareto.py}.

\texttt{hypervolume\_2d} uses a reference CNOT count derived from the current
archive maximum.  This makes the value convenient as a rough within-run summary,
but it should not be compared across runs unless the reference point is fixed.

\subsection{Lower Priority: Dead or Partially Integrated Code Paths}

Examples:

\begin{itemize}[leftmargin=*]
  \item \texttt{trainer.py} defines \texttt{model\_forward} inside
  \texttt{\_compile\_jit\_fns} but does not use it.
  \item \texttt{loss.py:joint\_sequence\_log\_prob} supports joint discrete and
  continuous sequence log probabilities, but the training loop does not use it.
  \item \texttt{benchmark.py} contains unused JSONL/summary/metadata helpers.
\end{itemize}

These are not correctness failures by themselves, but they make it harder to see
the intended training objective and benchmark workflow.

\section{Bottlenecks to Address Before 4-Qubit Sweeps}

\begin{enumerate}[leftmargin=*]
  \item \textbf{Inner refinement on every sample.}  With
  \texttt{num\_samples=2048}, \texttt{inner\_refine\_steps=16}, and a 4-qubit
  horizon near \(458\), refinement is likely the dominant cost.

  \item \textbf{Transformer sequence length.}  PPO training uses full-sequence
  causal attention.  The jump from \(T=105\) at 3 qubits to \(T=459\) at 4
  qubits is large.

  \item \textbf{Dense exact evaluation.}  Dense \(16 \times 16\) matrices are
  still manageable for 4 qubits, but not when multiplied by long sequences, large
  batches, and many refinement steps.

  \item \textbf{Host/device transfer and Python loops.}  Each epoch converts
  rollout outputs to NumPy, pushes buffer entries in a Python loop, builds Python
  Pareto objects, and may run host-side sweep polishing.  These costs grow with
  sample count and archive size.

  \item \textbf{Qiskit reference compilation.}  Compiling dense
  \texttt{UnitaryGate} targets can become expensive for 4+ qubits and should be
  optional in training-oriented runs.

  \item \textbf{JAX compilation/cache setup.}  Current cache warnings indicate
  the persistent compilation cache may not be usable in some environments.  This
  is painful when changing qubit count, horizon, or batch size because JAX must
  recompile large functions.
\end{enumerate}

\section{Suggested 4-Qubit Roadmap}

\subsection{Phase 0: Correctness Before New Claims}

Before trusting 4-qubit results, fix or guard:

\begin{itemize}[leftmargin=*]
  \item one consistent angle-array convention across refinement, simplification,
  reporting, and artifacts;
  \item fidelity verification after simplification;
  \item tests for \(\mathrm{SX}^4\), \(R_P(2\pi)\) global-phase removal if that
  remains part of the objective, and report-circuit fidelity after
  simplification;
  \item artifact storage of token sequences and angles.
\end{itemize}

\subsection{Phase 1: Structured 4-Qubit Brickwork}

Start with brickwork targets, not full Haar targets.  Recommended experiment
shape:

\begin{itemize}[leftmargin=*]
  \item nearest-neighbor CNOT pool first, then compare all-to-all;
  \item horizon \(64\) to \(128\), increased only if needed;
  \item inner refinement on top-\(K\) samples rather than all samples;
  \item online simplification before reward and archive update;
  \item per-seed JSONL writes so long runs are resumable.
\end{itemize}

\subsection{Phase 2: Broader 4-Qubit Targets}

After structured targets are stable, try harder 4-qubit targets:

\begin{itemize}[leftmargin=*]
  \item random reachable targets with known generating circuits;
  \item deeper brickwork targets;
  \item dense Haar targets only with reduced sample counts and staged
  refinement.
\end{itemize}

Use Qiskit as a reference, but decouple Qiskit compilation from the training loop
so training runs do not fail or slow down because of reference synthesis.

\subsection{Phase 3: 5+ Qubits}

For 5+ qubits, direct dense-unitary sequence search should not be the main plan.
Use hierarchical decomposition, local-window GQE, tensor-network or sketch-based
screening, and exact dense verification only for final candidates where feasible.

\section{Bottom Line}

The current 2-qubit improvement is a useful signal: the policy is finding
high-fidelity circuits with substantial removable redundancy.  To make the
method scale, the simplifier should become part of the learning and archiving
loop, not only a final reporting step.  The most urgent correctness issue is the
angle-alignment mismatch after simplification, because it can make exported
Qiskit circuits disagree with archive fidelities.  The most urgent scalability
issue is the combination of long 4-qubit horizons and per-sample inner
refinement.

\end{document}
